% --- Appendix A: Formal Metric Definitions ---

\label{app:metrics}

This appendix collects formal definitions for the five quantitative
metrics in the main text and the bootstrap procedure used for all
reported confidence intervals.

\subsection{NMACR Aggregate}

For each (task, model) pair, the Path~B rubric produces five
dimension-level scores $N, M, A, C, R \in [0,1]$. The aggregate
from~(\ref{eq:nmacr}) is
\[
  S_{\text{NMACR}} = 0.10\,N + 0.20\,M + 0.30\,A + 0.15\,C + 0.25\,R.
\]
The ordering $A > R > M > C > N$ weights assumption articulation
highest in a Bayesian task and a correct number with no derivation
lowest. Per-model accuracy is the mean of $S_{\text{NMACR}}$ over
the 171 base tasks. The literature trail behind the weights lives
in the scoring metadata.

\subsection{Robustness Pass-Flip}

Equation~(\ref{eq:robustness}) defines per-model robustness as one
minus the mean absolute pass-flip between a base task and the mean
pass over its perturbations. Equivalently, for each base task with
$k$ perturbations, the contribution to $R_m$ is
\[
  \frac{1}{k}\sum_{p \in \mathcal{P}(t)}
    \mathbb{1}\bigl[\mathrm{pass}(t, m) = \mathrm{pass}(p, m)\bigr].
\]
The two formulations agree exactly for binary $\mathrm{pass}(\cdot)$.

\subsection{Expected Calibration Error}

Binary-pass ECE is computed with $B = 4$ confidence buckets parsed
from the response's hedging language (low, medium, high, unstated).
Following Guo et al.~\cite{guo2017calibration} and
restating~(\ref{eq:ece}),
\[
  \mathrm{ECE} = \sum_{b=1}^{4} \frac{|B_b|}{N}
                 \,\bigl|\,\mathrm{acc}(B_b) - \mathrm{conf}(B_b)\,\bigr|,
\]
where bucket centres for ``low'', ``medium'', ``high'', and
``unstated'' map to $0.25$, $0.50$, $0.85$, and $0.50$ before
computing $\mathrm{conf}(B_b)$.

\subsection{Krippendorff's \texorpdfstring{$\alpha$}{alpha}}

For each judged dimension, inter-rater reliability between the
keyword rubric and the external judge is
\begin{equation}
  \alpha = 1 - \frac{D_o}{D_e},
  \label{eq:krippendorff}
\end{equation}
with $D_o$ observed disagreement among raters and $D_e$ disagreement
expected under chance. Values near $1$ indicate strong agreement;
near $0$, chance. We compute $\alpha$ as defined in
Eq.~\ref{eq:krippendorff} on the three Path~B dimensions
(method-structure, assumption-compliance, reasoning-quality) used
to characterize keyword-vs-judge alignment in
Section~\ref{sec:results}.

\subsection{Spearman Rank Correlation}

Per-dimension keyword-vs-judge agreement is also reported as
\begin{equation}
  \rho = 1 - \frac{6 \sum_{i=1}^{n} d_i^2}{n(n^2 - 1)},
  \label{eq:spearman}
\end{equation}
with $d_i$ the rank difference between the keyword-rubric and
external-judge score for the $i$-th response in the paired sample
of size $n$. We report $\rho$ as defined in Eq.~\ref{eq:spearman}
rather than Pearson $r$ because the rubric-judge relationship is
monotonic but not linear.

\subsection{Bootstrap Procedure}

All reported confidence intervals use the 95\% percentile bootstrap
on per-run aggregates, with $B = 10{,}000$ resamples and a fixed
RNG seed (42). The resampling unit is the individual (task, model)
run; resamples preserve the per-model count ($n = 171$ for accuracy,
$n = 398$ for the v2 perturbation pool used in robustness). CIs are
reported throughout the results section as $[\text{lower},
\text{upper}]$ at the $0.025$ and $0.975$ quantiles of the
resampled statistic distribution.
